% ═══════════════════════════════════════════════════════════════
% MUSTERLÖSUNG: Die Gewichtskraft (UE 3)
% Physik Klasse 7 – Gesamtschule MV – 45 Minuten
% Basiert auf: MINT Arbeitsblatt Template v1.0
% ═══════════════════════════════════════════════════════════════

\documentclass[11pt, a4paper]{article}

\usepackage[utf8]{inputenc}
\usepackage[T1]{fontenc}
\usepackage[ngerman]{babel}

% --- SCHRIFTEN ---
\usepackage{helvet}
\renewcommand{\familydefault}{\sfdefault}

\usepackage{microtype}
\usepackage[margin=1.5cm, top=2cm, bottom=1.5cm]{geometry}
\usepackage{enumitem}
\usepackage{tabularx}
\usepackage{booktabs}
\usepackage{array}
\usepackage{tikz}
\usetikzlibrary{calc}

\usepackage{amsmath, amssymb}
\usepackage{siunitx}
\sisetup{locale = DE, per-mode = symbol}

\usepackage{fancyhdr}
\usepackage{lastpage}
\usepackage{needspace}

\usepackage{xcolor}
\usepackage{tcolorbox}
\tcbuselibrary{skins}

% ═══════════════════════════════════════════════════════════════
% FARBEN
% ═══════════════════════════════════════════════════════════════

\definecolor{physik}{HTML}{2c5aa0}
\definecolor{hellgrau}{HTML}{F5F5F5}
\definecolor{mittelgrau}{HTML}{888888}
\definecolor{rahmen}{HTML}{CCCCCC}
\definecolor{loesung}{HTML}{c62828}

\colorlet{fachfarbe}{physik}

% ═══════════════════════════════════════════════════════════════
% VARIABLEN
% ═══════════════════════════════════════════════════════════════

\newcommand{\fach}{Physik}
\newcommand{\thema}{Die Gewichtskraft}
\newcommand{\klasse}{7}
\newcommand{\stunde}{03}
\newcommand{\maxpunkte}{26}

% Lösungsfarbe
\newcommand{\lsg}[1]{\textcolor{loesung}{\textbf{#1}}}

% ═══════════════════════════════════════════════════════════════
% KOPF- UND FUSSZEILE
% ═══════════════════════════════════════════════════════════════

\pagestyle{fancy}
\fancyhf{}
\renewcommand{\headrulewidth}{0.4pt}
\renewcommand{\footrulewidth}{0pt}
\lhead{\small \fach{} Klasse \klasse{} | \thema}
\chead{\small \textcolor{loesung}{\textbf{MUSTERLÖSUNG}}}
\rhead{\small Name: \underline{\hspace{3.5cm}}}
\cfoot{\small\textcolor{mittelgrau}{Seite \thepage{} von \pageref{LastPage}}}

% ═══════════════════════════════════════════════════════════════
% BOXEN
% ═══════════════════════════════════════════════════════════════

% Merkkasten
\newtcolorbox{merkkasten}[1][]{
    colback=hellgrau,
    colframe=hellgrau,
    coltitle=black,
    colbacktitle=hellgrau,
    boxrule=0pt,
    arc=0pt,
    left=8pt, right=8pt, top=8pt, bottom=6pt,
    borderline north={2pt}{0pt}{fachfarbe},
    before skip=6pt, after skip=6pt,
    fonttitle=\bfseries,
    #1
}

% Formelkasten
\newtcolorbox{formelkasten}{
    colback=white,
    colframe=fachfarbe,
    boxrule=1pt,
    arc=0pt,
    left=8pt, right=8pt, top=6pt, bottom=6pt,
    before skip=4pt, after skip=4pt
}

% Hinweiskasten
\newtcolorbox{hinweiskasten}{
    colback=white,
    colframe=white,
    boxrule=0pt,
    arc=0pt,
    left=8pt, right=4pt, top=4pt, bottom=4pt,
    borderline west={3pt}{0pt}{fachfarbe}
}

% ═══════════════════════════════════════════════════════════════
% BEFEHLE
% ═══════════════════════════════════════════════════════════════

% Differenzierung anzeigen - ausgeschaltet auf Wunsch
\newif\ifzeigesternchen
\zeigesternchenfalse

\newcommand{\punktebox}[1]{\hfill\small\textcolor{mittelgrau}{[\_\_/#1]}}

\newcommand{\stern}[1]{%
    \ifzeigesternchen%
        \textsuperscript{\textcolor{fachfarbe}{(#1)}}%
    \fi%
}

\newcommand{\aufgabe}[2]{%
    \needspace{4\baselineskip}%
    \vspace{10pt}%
    \noindent\textbf{\textcolor{fachfarbe}{Aufgabe #1}} #2%
    \vspace{4pt}%
    \nopagebreak[4]%
}

\newcommand{\operator}[1]{\textbf{\textcolor{fachfarbe}{#1}}}

\newlist{teil}{enumerate}{2}
\setlist[teil,1]{label=\alph*), leftmargin=1.8em, itemsep=3pt, parsep=0pt, topsep=3pt}
\setlist[teil,2]{label=\roman*), leftmargin=1.8em, itemsep=2pt}

\setlength{\parindent}{0pt}
\setlength{\parskip}{4pt}

% ═══════════════════════════════════════════════════════════════
% DOKUMENT
% ═══════════════════════════════════════════════════════════════

\begin{document}

% ═══════════════════════════════════════════════════════════════
% HEADER
% ═══════════════════════════════════════════════════════════════

\begin{center}
    {\Large\bfseries\textcolor{fachfarbe}{\thema}}\\[0.2em]
    {\small Arbeitsblatt | \fach{} Klasse \klasse{} | Kräfte}
\end{center}
\vspace{-0.3em}
\hrule
\vspace{0.6em}

% ═══════════════════════════════════════════════════════════════
% MERKKASTEN 1: MASSE
% ═══════════════════════════════════════════════════════════════

\begin{merkkasten}[title={Merksatz 1: Die Masse}]
\punktebox{3}

\vspace{-0.5em}
Die \textbf{Masse} $m$ gibt an, wie viel \lsg{Stoff} in einem Körper enthalten ist.

\vspace{0.3em}
\begin{tabular}{@{}ll@{}}
\textbf{Formelzeichen:} & $m$ \\
\textbf{Einheit:} & \lsg{kg} (Kilogramm) \\
\textbf{Wichtig:} & Die Masse ist \lsg{überall gleich}, das bedeutet, \\
                  & sie ändert sich nicht, egal wo sich der Körper befindet!
\end{tabular}
\end{merkkasten}

% ═══════════════════════════════════════════════════════════════
% MERKKASTEN 2: GEWICHTSKRAFT
% ═══════════════════════════════════════════════════════════════

\begin{merkkasten}[title={Merksatz 2: Die Gewichtskraft}]
\punktebox{3}

\vspace{-0.5em}
\begin{minipage}[t]{0.62\textwidth}
Die \textbf{Gewichtskraft} $F_G$ ist die Kraft, mit der ein Himmelskörper (z.\,B. die Erde) einen Körper \lsg{anzieht}.

\vspace{0.5em}
\begin{formelkasten}
\centering
{\Large $F_G = m \cdot g$}
\end{formelkasten}

\vspace{0.3em}
\begin{tabular}{@{}r@{ = }l@{}}
$F_G$ & Gewichtskraft in \lsg{N (Newton)} \\
$m$   & Masse in \si{\kilogram} \\
$g$   & Ortsfaktor in \lsg{N/kg}
\end{tabular}
\end{minipage}
\hfill
\begin{minipage}[t]{0.32\textwidth}
\centering
\begin{tikzpicture}[scale=0.9]
    % Masse (Rechteck)
    \fill[fachfarbe!20] (-0.7,0) rectangle (0.7,0.9);
    \draw[fachfarbe, thick] (-0.7,0) rectangle (0.7,0.9);
    \node at (0,0.45) {$m$};
    % Kraftpfeil nach unten
    \draw[->, line width=2pt, red!70!black] (0,0) -- (0,-1.6);
    \node[red!70!black, right] at (0.15,-0.8) {$F_G$};
    % Aufhängung
    \draw[thick] (0,0.9) -- (0,1.3);
    \fill[black] (0,1.3) circle (0.08);
\end{tikzpicture}

\vspace{0.3em}
\textbf{Faustformel:}\\
\SI{100}{\gram} $\hat{=}$ \SI{1}{\newton}
\end{minipage}

\end{merkkasten}

% ═══════════════════════════════════════════════════════════════
% HINWEIS: ORTSFAKTOR
% ═══════════════════════════════════════════════════════════════

\begin{hinweiskasten}
\textbf{Der Ortsfaktor $g$} wird auch \textbf{Fallbeschleunigung} oder \textbf{Gravitationsfeldstärke} genannt.

\vspace{0.3em}
Auf der Erde gilt genau: $g = \SI{9,81}{\newton\per\kilogram}$

\vspace{0.2em}
\textbf{Wir rechnen näherungsweise mit:} \quad {\large $g = \SI{10}{\newton\per\kilogram}$} \quad (kopfrechenbar!)

\vspace{0.3em}
Der Ortsfaktor ist \textbf{ortsabhängig} -- auf dem Mond oder anderen Planeten ist $g$ anders!
\end{hinweiskasten}

\vspace{0.3em}

% ═══════════════════════════════════════════════════════════════
% AUFGABEN
% ═══════════════════════════════════════════════════════════════

\aufgabe{1}{\stern{*}} \punktebox{4}

\operator{Berechne} die Gewichtskraft $F_G$ auf der Erde ($g = \SI{10}{\newton\per\kilogram}$).

\vspace{0.2em}
\begin{teil}
    \item $m = \SI{5}{\kilogram}$ \hfill $F_G = $ \lsg{\SI{5}{\kilogram} $\cdot$ \SI{10}{\newton\per\kilogram}} $= $ \lsg{\SI{50}{\newton}}
    \item $m = \SI{8}{\kilogram}$ \hfill $F_G = $ \lsg{\SI{8}{\kilogram} $\cdot$ \SI{10}{\newton\per\kilogram}} $= $ \lsg{\SI{80}{\newton}}
\end{teil}

% ---

\aufgabe{2}{\stern{*}} \punktebox{3}

\operator{Rechne um} mit der Faustformel: \SI{100}{\gram} $\hat{=}$ \SI{1}{\newton}

\vspace{0.2em}
\begin{teil}
    \item $m = \SI{200}{\gram}$ \hfill $F_G = $ \lsg{\SI{2}{\newton}}
    \item $m = \SI{500}{\gram}$ \hfill $F_G = $ \lsg{\SI{5}{\newton}}
    \item $m = \SI{1}{\kilogram} = \SI{1000}{\gram}$ \hfill $F_G = $ \lsg{\SI{10}{\newton}}
\end{teil}

% ---

\aufgabe{3}{\stern{**}} \punktebox{4}

\operator{Berechne} die Gewichtskraft. Achte auf die Einheiten!

\vspace{0.2em}
\begin{teil}
    \item $m = \SI{2,5}{\kilogram}$ \hfill $F_G = $ \lsg{\SI{2,5}{\kilogram} $\cdot$ \SI{10}{\newton\per\kilogram}} $= $ \lsg{\SI{25}{\newton}}
    \item $m = \SI{0,4}{\kilogram}$ \hfill $F_G = $ \lsg{\SI{0,4}{\kilogram} $\cdot$ \SI{10}{\newton\per\kilogram}} $= $ \lsg{\SI{4}{\newton}}
\end{teil}

% ═══════════════════════════════════════════════════════════════
% SEITE 2
% ═══════════════════════════════════════════════════════════════
\newpage

\aufgabe{4}{\stern{**}} \punktebox{4}

\operator{Berechne} die Masse. Stelle dazu die Formel um: $m = \dfrac{F_G}{g}$

\vspace{0.2em}
\begin{teil}
    \item $F_G = \SI{40}{\newton}$ \hfill $m = $ \lsg{$\dfrac{\SI{40}{\newton}}{\SI{10}{\newton\per\kilogram}}$} $= $ \lsg{\SI{4}{\kilogram}}
    \item $F_G = \SI{150}{\newton}$ \hfill $m = $ \lsg{$\dfrac{\SI{150}{\newton}}{\SI{10}{\newton\per\kilogram}}$} $= $ \lsg{\SI{15}{\kilogram}}
\end{teil}

% ---

\aufgabe{5}{\stern{***}} \punktebox{3}

Der Ortsfaktor $g$ ist nicht überall gleich! \operator{Berechne} die Gewichtskraft eines Astronauten mit Anzug ($m = \SI{100}{\kilogram}$) auf verschiedenen Himmelskörpern.

\vspace{0.3em}
\begin{center}
\begin{tabular}{|l|c|c|}
\hline
\textbf{Ort} & \textbf{Ortsfaktor $g$} & \textbf{Gewichtskraft $F_G$} \\
\hline
Erde & \SI{10}{\newton\per\kilogram} & \lsg{\SI{1000}{\newton}} \\
\hline
Mond & \SI{1,6}{\newton\per\kilogram} & \lsg{\SI{160}{\newton}} \\
\hline
Mars & \SI{4}{\newton\per\kilogram} & \lsg{\SI{400}{\newton}} \\
\hline
\end{tabular}
\end{center}

% ---

\aufgabe{6}{\stern{***}} \punktebox{2}

\operator{Erkläre} den Unterschied zwischen Masse und Gewichtskraft.

\vspace{0.2em}
\begin{teil}
    \item Warum bleibt die \textbf{Masse} eines Körpers überall gleich?

    \lsg{Die Masse gibt an, wie viel Stoff in einem Körper enthalten ist. Diese Stoffmenge ändert sich nicht, egal wo sich der Körper befindet.}

    \item Warum ändert sich die \textbf{Gewichtskraft} je nach Ort?

    \lsg{Die Gewichtskraft hängt vom Ortsfaktor g ab. Der Ortsfaktor ist auf verschiedenen Himmelskörpern unterschiedlich groß, weil sie unterschiedliche Massen haben.}
\end{teil}

% ═══════════════════════════════════════════════════════════════
% PUNKTE
% ═══════════════════════════════════════════════════════════════

\vfill
\hrule
\vspace{0.4em}

\hfill\textbf{Punkte:} \underline{\hspace{1.2cm}} / \maxpunkte

\end{document}
